\begin{proposition}[Refuted printed surjectivity and positivity, with a second example]\hypertarget{tag-01FY}{}\label{cdp:st:D.1.13}
\StatementRecord{cdp:st:D.1.13}{proved}{01FY}{\SeriesRef{cdp:st:D.1.6}, \SeriesRef{cdp:st:D.1.12}, \SeriesRef{cdp:st:D.1.9}, \SeriesRef{cdp:st:D.1.4}, \SeriesRef{cdp:st:E1}}


For Wang's threefold $Y=X_2(\lambda)$ in \ref{cdp:st:D.1.6} one has
\[
K_Y\simeq f^*\OO_{\PP^1}(-1),\qquad
R^1f_*\OO_Y\simeq\OO\oplus\OO(-1),\qquad
R^2f_*\OO_Y\simeq\OO(-1).
\]
Consequently $h^q(Y,\OO_Y)=(1,1,0,0)$, and restriction to every
smooth two-torus fibre has rank one. This threefold satisfies
$H^1(Y;\Z)=H^2(Y;\Z)=0$, $a(Y)=1$, and is a second
counterexample to CDP20 Corollary 5.2 and Proposition 5.5.
Its integral homology torsion distinguishes it from the retained
integral homology sphere.
\end{proposition}


\begin{proof}

\begin{enumerate}[label={$\langle1\rangle$\arabic*.},leftmargin=*]
\item We compute the canonical line through the actual local fillings.
The rational elliptic surface $\pi:S\to\PP^1$ in this construction
has Weierstrass equation
\[
y^2=x^3-3(1+t)x+(2+3t+2t^2)
\]
and fundamental line bundle $\OO(1)$; hence
$K_S=\pi^*\OO(-1)$. On the complement of the zero section in
$M=\OO_S(P-O)$, the determinant of the relative cotangent sequence
identifies its canonical line with the pullback of $K_S$.
In a local multiplicative coordinate $v$, this identification sends
a local two-form $\omega_S$ to $\omega_S\wedge d\log v$.
Changing $v$ by a nonzero holomorphic function on $S$ adds a
one-form from $S$ to $d\log v$, whose wedge with $\omega_S$ is zero.
This verifies the coordinate independence of the identification.
The contracting generator covers translation by $-6P$ on the
elliptic fibres. Translation preserves the relative elliptic
differential, and the extra base terms disappear after wedging
with the base differential. Its line-bundle lift changes $d\log v$
only by a base one-form. It therefore preserves this canonical
identification. The latter descends to
$K_Y=f^*\OO(-1)$ on the open quotient.

\item At either rank-one filling, let $q$ be the base coordinate and
$(x,y)$ the two multiplicative coordinates of the construction.
At $t=0$ the explicit change is $x=v$, $y=v^2/z$; its exponent
matrix is unimodular. At $\alpha_+$ no change is required.
Thus the preceding three-form is a nonzero holomorphic unit times
\[
dq\wedge d\log x\wedge d\log y.
\]
The nondegenerating $y$ direction becomes an elliptic curve after
quotienting by its contracting period. Its logarithmic differential
is a nonvanishing relative one-form; changes by a base differential
do not affect the displayed wedge. At a double curve the equation
is $uv=q$, and
$dq\wedge d\log u=-du\wedge dv$. Therefore the three-form extends
without a zero or pole over the full rank-one filling.

\item At $\alpha_-$ the rank-two filling uses a unimodular period lattice
and the standard toric charts $q=z_0z_1z_2$. A pair of independent
fibre characters, together with $q$, is a unimodular system of
characters in each smooth chart. The same form is, up to a nonzero
constant and a holomorphic unit,
\[
q\,d\log z_0\wedge d\log z_1\wedge d\log z_2
=dz_0\wedge dz_1\wedge dz_2.
\]
It extends without zero or pole, including at the triple points.

\item It remains to check the order-three logarithmic filling at infinity.
Put $s=1/t$, $x_s=s^2x$ and $y_s=s^3y$. Its Weierstrass
coefficients have orders three and four in $s$, giving type $IV^*$.
On the cubic base change $s=r^3$, put
$x_s=r^4X$, $y_s=r^6Y$. The good-reduction equation is
\[
Y^2=X^3-3r(r^3+1)X+(2+3r^3+2r^6),
\]
whose central elliptic curve is smooth. The canonical form satisfies
\[
ds\wedge\frac{dx_s}{y_s}
=3\,dr\wedge\frac{dX}{Y}.
\]
It has no zero or pole. For $\zeta=e^{2\pi i/3}$ the root deck
action has weights $r\mapsto\zeta r$, $X\mapsto\zeta^2X$,
$Y\mapsto Y$; the weights on $dr$ and $dX/Y$ multiply to one.
The additional invariant translation used in the logarithmic
transform preserves the holomorphic two-form on the torus cover.
Hence the three-form is invariant and descends through the free
cyclic action as a nonvanishing three-form. This proves that the
canonical identification from the open quotient extends everywhere,
giving $K_Y=f^*\OO(-1)$.

\item The map is flat and Gorenstein, since its source is smooth and its
base is a smooth curve. Relative Serre duality gives
\[
(R^2f_*\OO_Y)^*=f_*\omega_{Y/\PP^1}
=f_*f^*\OO(1)=\OO(1).
\]
Statement \ref{cdp:st:D.1.12} gives local freeness and
$\det(R^1f_*\OO_Y)=R^2f_*\OO_Y=\OO(-1)$.
The nowhere-zero section from \ref{cdp:st:D.1.9} embeds $\OO$ as a subbundle
of $R^1f_*\OO_Y$ with quotient $\OO(-1)$. The extension splits
because $H^1(\PP^1,\OO(1))=0$. Coherent Leray now computes all
four displayed structure-sheaf cohomology groups, and \ref{cdp:st:D.1.4}
identifies the degree-one restriction with rank-one evaluation.

\item The homology computation in \ref{cdp:st:D.1.6} and the universal coefficient
theorem give $H^1(Y;\Z)=H^2(Y;\Z)=0$; specifically the torsion
in $H_2$ contributes to $H^3$, not to $H^2$.
The general fibre is K\"ahler and $b_2(Y)=0$, so \ref{cdp:st:E1} proves that
this map is the algebraic reduction and $a(Y)=1$.
The retained six-sphere has zero second integral homology, whereas
$H_2(Y;\Z)=\Z/2\Z$. An isomorphism of complex manifolds would
induce an isomorphism of these groups, which is impossible.
This completes the independent counterexample.

\end{enumerate}
\end{proof}